\subsection{\textit{Enforcement} del programa de inversiones}

Excluir el avance físico o financiero del ICG no implica reducir la exigibilidad del programa de inversiones. Las inversiones reconocidas en la tarifa constituyen obligaciones regulatorias cuya ejecución debe permanecer sujeta a programación, seguimiento, fiscalización y sanción cuando el incumplimiento sea atribuible a la empresa prestadora. Como se mostró en la Sección~4.3, este mecanismo puede preservar el incentivo hacia la ejecución siempre que la pérdida esperada por incumplimiento sea suficientemente sensible al avance y tenga efectos observables sobre la reputación de la empresa y la evaluación de sus administradores.

El Programa de Ejecución de Metas Físico-Financieras puede cumplir una función central en este esquema. Para ello, debe permitir identificar las inversiones y medidas de mejora, sus fuentes de financiamiento, los hitos de ejecución, los plazos previstos, los riesgos relevantes y las medidas de contingencia. Esta trazabilidad facilita verificar si la empresa adoptó oportunamente las acciones necesarias y si los retrasos observados se encontraban bajo su control.

La efectividad del mecanismo depende también de la oportunidad de la fiscalización. El control debería combinar evaluaciones periódicas, alertas tempranas y medidas correctivas con consecuencias sancionadoras frente a incumplimientos atribuibles. Ello permite distinguir entre un desvío temporal que todavía puede corregirse y una postergación efectiva de la inversión, considerando la fase del ciclo de inversión, la materialidad del retraso y las acciones adoptadas por la empresa.

La imputabilidad constituye un elemento central. La empresa no debería ser sancionada por factores que no podía razonablemente controlar, siempre que haya identificado los riesgos y adoptado las acciones de gestión y contingencia exigibles. Sin embargo, las restricciones administrativas, contractuales o institucionales no deberían operar como justificaciones automáticas cuando la empresa pudo anticiparlas, prevenirlas o mitigar sus efectos.

Finalmente, el mecanismo debe ser creíble y visible. Una sanción tardía, poco difundida o desconectada de la administración responsable puede generar incentivos débiles, incluso cuando su monto formal sea elevado. La publicación de resultados, la comparación entre empresas y la identificación de responsabilidades institucionales complementan el efecto monetario de la sanción. La condición bajo la cual este \textit{enforcement} separado permite excluir el avance físico del ICG sin reducir la calidad esperada se desarrolla en la Sección~4.4.
